\chapter{Iloczyn Cauchy'ego szeregów}
\label{ma1:ch:cauchy-product}

Mnożenie dwóch skończonych sum prowadzi do sumy iloczynów wszystkich par składników. Dla szeregów nieskończonych pojawia się problem: w jakiej kolejności należy dodawać te iloczyny i kiedy otrzymana suma jest iloczynem sum wyjściowych? Rozstrzygniemy to dla szeregów rzeczywistych i zespolonych, nie korzystając jeszcze z teorii funkcji wykładniczej ani szeregów potęgowych.

\section{Definicja i geometria indeksów}
\begin{definicja}
\label{ma1:def:cauchy-product}
Niech $K=\R$ lub $K=\C$, a $(a_n)_{n\ge0}$ i $(b_n)_{n\ge0}$ będą ciągami w $K$. Iloczynem Cauchy'ego szeregów $\sum_{n=0}^\infty a_n$ i $\sum_{n=0}^\infty b_n$ nazywamy szereg $\sum_{n=0}^\infty c_n$, gdzie dla każdego całkowitego $n\ge0$
\[
 c_n=\sum_{j=0}^n a_jb_{n-j}.
\]
Definicja nie wymaga zbieżności żadnego z trzech szeregów.
\end{definicja}
Indeksowanie od zera upraszcza wzory: wyraz $c_n$ zawiera dokładnie te iloczyny $a_jb_k$, dla których $j+k=n$. Dla szeregów indeksowanych od jedynki można najpierw przeindeksować wyrazy, przyjmując $\widetilde a_n=a_{n+1}$ i $\widetilde b_n=b_{n+1}$.

\begin{lemat}[Tożsamości skończone]
\label{ma1:lem:cauchy-finite-identities}
Niech $(a_n),(b_n)$ będą ciągami w $K=\R$ lub $K=\C$, indeksowanymi od zera, i niech $(c_n)$ będzie określony jak w definicji \ref{ma1:def:cauchy-product}. Dla każdego całkowitego $N\ge0$ oznaczmy
$A_N=\sum_{j=0}^N a_j$, $B_N=\sum_{k=0}^N b_k$ i $C_N=\sum_{n=0}^N c_n$. Wówczas
\[
 C_N=\sum_{\substack{j,k\ge0\\j+k\le N}}a_jb_k
 =\sum_{j=0}^N a_jB_{N-j},
 \qquad
 A_NB_N=\sum_{j=0}^N\sum_{k=0}^N a_jb_k.
\]
We wszystkich sumach indeksy są całkowite.
\end{lemat}
\begin{proof}
Rozwijamy $C_N=\sum_{n=0}^N\sum_{j=0}^n a_jb_{n-j}$. Przypisanie $(n,j)\mapsto(j,n-j)$ jest bijekcją zbioru par całkowitych $0\le j\le n\le N$ na zbiór par nieujemnych $j,k$ z $j+k\le N$. Przypisaniem odwrotnym jest $(j,k)\mapsto(j+k,j)$. Dla ustalonego $j$ dopuszczalne wartości $k$ to $0,\ldots,N-j$, co daje drugą równość. Ostatnia równość wynika z rozdzielności mnożenia względem dodawania. Wszystkie przekształcenia dotyczą sum skończonych.
\end{proof}

\begin{figure}[htbp]
\centering
\begin{tikzpicture}[x=.8cm,y=.8cm,>=Stealth]
\fill[blue!10] (-.22,-.22) rectangle (4.22,4.22);
\fill[teal!22] (-.22,-.22)--(4.22,-.22)--(4.22,.22)--(.22,4.22)--(-.22,4.22)--cycle;
\draw[->] (-.4,0)--(5.3,0) node[right] {$j$};
\draw[->] (0,-.4)--(0,5.3) node[above] {$k$};
\foreach \j in {0,...,4}{
 \foreach \k in {0,...,4}{
  \pgfmathtruncatemacro{\d}{\j+\k}
  \ifnum\d>4
   \fill[blue!65!black] (\j,\k) circle (2pt);
  \else
   \fill[teal!75!black] (\j,\k) circle (2pt);
  \fi
 }
}
\draw[teal!75!black,thick] (0,4)--(4,0);
\draw[densely dashed,blue!65!black] (0,4)--(4,4)--(4,0);
\node[left] at (0,4) {$N$};
\node[below] at (4,0) {$N$};
\node[below left] at (0,0) {$0$};
\node[anchor=west,align=left] at (5.5,3.6) {\small Trójkąt: $j+k\le N$\\\small suma $C_N$};
\node[anchor=west,align=left] at (5.5,1.6) {\small Poza trójkątem: $j+k>N$\\\small część kwadratu $j,k\le N$};
\fill[teal!75!black] (5.25,3.65) circle (2pt);
\fill[blue!65!black] (5.25,1.65) circle (2pt);
\end{tikzpicture}
\caption{Dla $N=4$ suma częściowa iloczynu Cauchy'ego obejmuje zielone punkty trójkąta, a iloczyn sum częściowych wszystkie zielone i niebieskie punkty kwadratu. Odcinki i tło służą jedynie wyróżnieniu skończonych zbiorów indeksów.}
\label{ma1:fig:cauchy-index-sets}
\end{figure}
\FloatBarrier

\section{Iloczyn szeregów bezwzględnie zbieżnych}
\begin{twierdzenie}
\label{ma1:thm:absolute-cauchy-product}
Niech $(a_n),(b_n)$ będą ciągami w $\R$ lub $\C$, indeksowanymi od zera, i niech oba szeregi $\sum_{n=0}^\infty a_n$, $\sum_{n=0}^\infty b_n$ będą bezwzględnie zbieżne do odpowiednio $A,B$. Jeżeli $c_n=\sum_{j=0}^n a_jb_{n-j}$ dla $n\ge0$, to $\sum_{n=0}^\infty c_n$ jest bezwzględnie zbieżny, ma sumę $AB$ oraz
\[
 \sum_{n=0}^\infty |c_n|
 \le \left(\sum_{n=0}^\infty |a_n|\right)
      \left(\sum_{n=0}^\infty |b_n|\right).
\]
\end{twierdzenie}
\begin{proof}
Oznaczmy $\alpha=\sum_{n=0}^\infty|a_n|$, $\beta=\sum_{n=0}^\infty|b_n|$. Dla każdego całkowitego $N\ge0$ z nierówności trójkąta, lematu \ref{ma1:lem:cauchy-finite-identities} zastosowanego do modułów oraz nieujemności składników mamy
\[
 \sum_{n=0}^N|c_n|\le
 \sum_{\substack{j,k\ge0\\j+k\le N}}|a_j||b_k|
 \le \sum_{j=0}^N\sum_{k=0}^N |a_j||b_k|
 \le\alpha\beta.
\]
Sumy częściowe szeregu modułów są niemalejące i ograniczone z góry przez $\alpha\beta$. Z twierdzenia \ref{ma1:thm:positive-series} szereg modułów jest zbieżny, a jego suma spełnia podane oszacowanie.

Niech $A_N,B_N,C_N$ oznaczają sumy częściowe z lematu \ref{ma1:lem:cauchy-finite-identities}. Dla $m=\lfloor N/2\rfloor$, jeżeli $0\le j,k\le N$ i $j+k>N$, to $j>m$ lub $k>m$: w przeciwnym razie $j+k\le2m\le N$. Zatem
\[
 |A_NB_N-C_N|
 \le \beta\sum_{j>m}|a_j|+
      \alpha\sum_{k>m}|b_k|.
\]
Istotnie, sumę modułów po zbiorze brakujących par ograniczamy przez sumę po dwóch zbiorach: $m<j\le N$, $0\le k\le N$ oraz $0\le j\le N$, $m<k\le N$. Wspólne pary można policzyć dwukrotnie, gdyż składniki są nieujemne. Sumy po tych zbiorach rozkładają się na iloczyny i są nie większe od podanych oszacowań. Gdy $N\to\infty$, mamy $m\to\infty$, a obie reszty szeregów modułów dążą do zera, ponieważ $\sum_{j>m}|a_j|=\alpha-\sum_{j=0}^{m}|a_j|\to0$, a dla drugiego szeregu zachodzi analogiczna równość. Ponieważ $A_NB_N\to AB$, z rachunku granic otrzymujemy $C_N\to AB$.
\end{proof}

\begin{przyklad}
\label{ma1:ex:geometric-product}
Niech $q\in\C$ i $|q|<1$. Biorąc $a_n=b_n=q^n$ dla $n\ge0$, otrzymujemy $c_n=\sum_{j=0}^n q^jq^{n-j}=(n+1)q^n$ (dla $q=0$ stosujemy konwencję $q^0=1$). Oba szeregi geometryczne są bezwzględnie zbieżne, więc
\[
 \sum_{n=0}^\infty (n+1)q^n=\frac1{(1-q)^2}.
\]
Dla $q=1/2$ suma wynosi $4$. Wzór wynika z mnożenia szeregów, bez różniczkowania i bez funkcji potęgowych o wykładnikach rzeczywistych.
\end{przyklad}

\section{Twierdzenie Mertensa}
\begin{twierdzenie}[Mertensa]
\label{ma1:thm:mertens-product}
Niech $(a_n),(b_n)$ będą ciągami w $\R$ lub $\C$, indeksowanymi od zera. Załóżmy, że $\sum_{n=0}^\infty a_n=A$, $\sum_{n=0}^\infty b_n=B$ i co najmniej jeden z tych szeregów jest zbieżny bezwzględnie. Wtedy iloczyn Cauchy'ego, czyli szereg o wyrazach $c_n=\sum_{j=0}^n a_jb_{n-j}$, jest zbieżny do $AB$.
\end{twierdzenie}
\begin{proof}
Możemy założyć, że szereg $\sum a_n$ jest bezwzględnie zbieżny, ponieważ zamiana $j$ na $n-j$ pokazuje, że zamiana ról $(a_n),(b_n)$ nie zmienia $c_n$. Oznaczmy $\alpha=\sum_{n=0}^\infty|a_n|$ oraz dla każdego całkowitego $N\ge0$ przyjmijmy $A_N=\sum_{j=0}^N a_j$, $B_N=\sum_{k=0}^N b_k$, $C_N=\sum_{n=0}^N c_n$. Z $B_n\to B$ wynika ograniczoność ciągu $(B_n-B)$, więc istnieje $M\ge1$ takie, że $|B_n-B|\le M$ dla każdego $n\ge0$.

Z lematu \ref{ma1:lem:cauchy-finite-identities} otrzymujemy
\[
 C_N-A_N B=\sum_{j=0}^N a_j(B_{N-j}-B).
\]
Ustalmy $\varepsilon>0$. Wybierzmy $J\ge0$ z $M\sum_{j>J}|a_j|<\varepsilon/2$. Ponieważ $B_n\to B$, istnieje $L\ge0$ takie, że $|B_n-B|<\varepsilon/(2(1+\alpha))$ dla $n\ge L$. Dla $N\ge J+L$ oraz $0\le j\le J$ mamy $N-j\ge L$, zatem
\[
\begin{aligned}
 |C_N-A_N B|
 &\le\sum_{j=0}^J |a_j||B_{N-j}-B|
      +\sum_{j=J+1}^N |a_j||B_{N-j}-B|\\
 &\le\frac{\varepsilon\alpha}{2(1+\alpha)}
        +M\sum_{j>J}|a_j|
 <\varepsilon.
\end{aligned}
\]
Druga suma jest zerowa, jeśli $N=J$. Wobec tego $C_N-A_NB\to0$. Ponieważ $A_N\to A$, otrzymujemy $C_N\to AB$.
\end{proof}

\begin{przyklad}
Niech $a_n=2^{-n}$ oraz $b_n=(-1)^n/(n+1)$ dla $n\ge0$. Pierwszy szereg ma sumę $2$ i jest zbieżny bezwzględnie, drugi jest zbieżny warunkowo z kryterium Leibniza. Jeżeli $S=\sum_{n=0}^\infty(-1)^n/(n+1)$, twierdzenie Mertensa daje
\[
 \sum_{n=0}^\infty\left(\sum_{j=0}^n\frac{(-1)^{n-j}}{2^j(n-j+1)}\right)=2S.
\]
Iloczyn ten jest zbieżny warunkowo. Istotnie, z definicji $c_0=b_0$ i dla $n\ge1$ zachodzi $c_n=b_n+\frac12c_{n-1}$ (wydzielamy składnik $j=0$ i w pozostałej sumie zmieniamy indeks na $j-1$). Gdyby $\sum_{n=0}^\infty|c_n|$ był zbieżny, to szereg $\sum_{n=1}^\infty(|c_n|+\frac12|c_{n-1}|)$ również byłby zbieżny na mocy liniowości i niezmienniczości zbieżności po pominięciu skończonego początku. Z $|b_n|\le |c_n|+\frac12|c_{n-1}|$ dla $n\ge1$ i kryterium porównawczego wynikałaby zbieżność $\sum_{n=0}^\infty|b_n|$, sprzecznie z rozbieżnością szeregu harmonicznego.
\end{przyklad}

\section{Iloczyn dwóch szeregów zbieżnych warunkowo}
\begin{przyklad}
\label{ma1:ex:divergent-cauchy-product}
Niech $a_n=b_n=(-1)^n/\sqrt{n+1}$ dla całkowitych $n\ge0$. Oba szeregi są zbieżne z kryterium Leibniza: $1/\sqrt{n+1}$ jest ciągiem malejącym, a dla $\varepsilon>0$ nierówność $n+1>1/\varepsilon^2$ daje $1/\sqrt{n+1}<\varepsilon$. Nie są zbieżne bezwzględnie, ponieważ $1/\sqrt{n+1}\ge1/(n+1)$, a szereg harmoniczny jest rozbieżny.
Jednak
\[
 c_n=(-1)^n\sum_{j=0}^n
       \frac1{\sqrt{(j+1)(n-j+1)}}.
\]
Dla $0\le j\le n$ mamy $j+1\le n+1$ i $n-j+1\le n+1$, zatem $\sqrt{(j+1)(n-j+1)}\le n+1$. Wszystkie składniki sumy bez znaku są dodatnie, a jest ich $n+1$. Multiplikatywność nieujemnych pierwiastków kwadratowych wynika z jednoznaczności pierwiastka: iloczyn dwóch nieujemnych pierwiastków jest nieujemny, a jego kwadrat jest iloczynem liczb podpierwiastkowych. Wobec tego
\[
 |c_n|\ge(n+1)\frac1{n+1}=1.
\]
Wyrazy iloczynu Cauchy'ego nie dążą do zera. Iloczyn ten jest rozbieżny na mocy twierdzenia \ref{ma1:thm:series-term-test}.
\end{przyklad}
Pytanie o istnienie sumy iloczynu Cauchy'ego trzeba odróżnić od pytania o jej wartość. Następne twierdzenie pokazuje, że jeśli wszystkie trzy szeregi są zbieżne, to suma iloczynu musi być iloczynem sum, nawet bez zbieżności bezwzględnej.

\section{Twierdzenie Abela o iloczynie Cauchy'ego}
\begin{lemat}
\label{ma1:lem:triangular-means}
Niech $(x_n)_{n\ge0},(y_n)_{n\ge0}$ będą ciągami w $\R$ lub $\C$ zbieżnymi odpowiednio do $x,y$. Wówczas
\[
 \lim_{N\to\infty}\frac1{N+1}\sum_{j=0}^N x_jy_{N-j}=xy.
\]
\end{lemat}
\begin{proof}
Ciąg $(y_n)$ jest ograniczony, więc wybieramy $M\ge1$ z $|y_n|\le M$ dla każdego $n$ oraz $|x|\le M$. Dla dowolnych całkowitych $j,k\ge0$ rozkład $x_jy_k-xy=(x_j-x)y_k+x(y_k-y)$ i nierówność trójkąta dają
$|x_jy_k-xy|\le M(|x_j-x|+|y_k-y|)$.
Ustalmy $\varepsilon>0$. Istnieje całkowite $J\ge1$ takie, że $|x_j-x|<\varepsilon/(4M)$ i $|y_j-y|<\varepsilon/(4M)$ dla każdego $j\ge J$. Jeżeli $N\ge2J$, to dla $J\le j\le N-J$ zachodzi $|x_jy_{N-j}-xy|<\varepsilon/2$.

Z ograniczoności obu ciągów istnieje rzeczywiste $D>0$ takie, że $|x_jy_k-xy|\le D$ dla wszystkich $j,k\ge0$. Na przykład dla ograniczeń $|x_j|\le P$, $|y_k|\le Q$ można wziąć $D=PQ+|xy|+1$. Dla $N\ge2J$ w sumie po $0\le j\le N$ dokładnie $2J$ indeksów nie spełnia $J\le j\le N-J$: po $J$ indeksów przy każdym końcu. W konsekwencji
\[
 \left|\frac1{N+1}\sum_{j=0}^N x_jy_{N-j}-xy\right|
 \le \frac{\varepsilon}{2}+\frac{2JD}{N+1}.
\]
Dla dostatecznie dużych $N$ drugi składnik jest mniejszy od $\varepsilon/2$, co dowodzi tezy.
\end{proof}
\begin{twierdzenie}[Abela o iloczynie Cauchy'ego]
\label{ma1:thm:abel-cauchy-product}
Niech $(a_n),(b_n)$ będą ciągami liczb rzeczywistych lub zespolonych, indeksowanymi od zera. Jeżeli szeregi $\sum_{n=0}^\infty a_n$, $\sum_{n=0}^\infty b_n$ oraz ich iloczyn Cauchy'ego $\sum_{n=0}^\infty c_n$, gdzie $c_n=\sum_{j=0}^n a_jb_{n-j}$, są zbieżne odpowiednio do $A,B,C$, to $C=AB$.
\end{twierdzenie}
\begin{proof}
Oznaczmy sumy częściowe przez $A_N,B_N,C_N$ jak w lemacie \ref{ma1:lem:cauchy-finite-identities}. Dla każdego całkowitego $N\ge0$ zachodzi tożsamość
\[
 \sum_{m=0}^N C_m=\sum_{j=0}^N A_j B_{N-j}.
\]
Aby ją uzasadnić, ustalmy parę nieujemnych indeksów całkowitych $r,s$ z $r+s\le N$. Składnik $a_rb_s$ po lewej występuje w sumach $C_m$ dokładnie dla $r+s\le m\le N$, a więc $N-r-s+1$ razy. Po prawej występuje dokładnie dla $r\le j\le N-s$, również $N-r-s+1$ razy. Jeżeli $r+s>N$, nie występuje po żadnej ze stron. Ponieważ rozpatrujemy skończone sumy, równość liczności wystąpień dla każdej pary dowodzi tożsamości.

Dzielimy obie strony przez $N+1$. Z lematu \ref{ma1:lem:triangular-means} zastosowanego do $x_j=A_j$, $y_j=B_j$ wynika, że prawa strona po podzieleniu dąży do $AB$. Stosując ten sam lemat do $x_j=C_j$ oraz stałego ciągu $y_j=1$, otrzymujemy, że lewa strona po podzieleniu dąży do $C$. Z jednoznaczności granicy mamy $C=AB$.
\end{proof}
Zaletą tego dowodu jest niezależność od teorii szeregów potęgowych. Wykorzystana własność średnich arytmetycznych pojawi się ponownie przy omawianiu metod sumowania szeregów.

\section{Sumowanie iloczynów po wszystkich parach indeksów}
\begin{twierdzenie}
\label{ma1:thm:pair-enumeration}
Niech $\N_0=\{0,1,2,\ldots\}$ i niech $(a_n)_{n\ge0},(b_n)_{n\ge0}$ będą ciągami w $\R$ lub $\C$ takimi, że oba szeregi są zbieżne bezwzględnie, odpowiednio do $A,B$. Dla dowolnej bijekcji $\pi:\N_+\to\N_0\times\N_0$, gdzie $\pi(l)=(j_l,k_l)$, szereg
\[
 \sum_{l=1}^\infty a_{j_l}b_{k_l}
\]
jest zbieżny bezwzględnie i ma sumę $AB$.
\end{twierdzenie}
\begin{proof}
Istnieje wyliczenie wszystkich par w kolejności rosnącej sumy współrzędnych: dla kolejnych $n=0,1,\ldots$ wypisujemy $(0,n),(1,n-1),\ldots,(n,0)$. Każda para $(j,k)$ należy do dokładnie jednego bloku, o numerze $n=j+k$, i występuje w nim dokładnie raz. Każdy blok jest skończony. Otrzymujemy bijekcję $\pi_0:\N_+\to\N_0\times\N_0$.
Oznaczmy $\alpha=\sum_{n\ge0}|a_n|$ i $\beta=\sum_{n\ge0}|b_n|$. Dla dowolnego skończonego początkowego fragmentu tego wyliczenia istnieje całkowite $M\ge0$ takie, że wszystkie jego pary mają obie współrzędne nie większe od $M$. Suma modułów odpowiadających tym parom jest nie większa od
$(\sum_{j=0}^M|a_j|)(\sum_{k=0}^M|b_k|)\le\alpha\beta$. Zatem szereg wyrazów $a_jb_k$ wyliczonych przez $\pi_0$ jest zbieżny bezwzględnie, a więc zbieżny.
Grupując jego kolejne bloki według $j+k=n$, otrzymujemy iloczyn Cauchy'ego. Twierdzenie \ref{ma1:thm:series-grouping} zapewnia zachowanie sumy, a twierdzenie \ref{ma1:thm:absolute-cauchy-product} identyfikuje ją z $AB$. Dowolna bijekcja $\pi$ wyznacza przestawienie wyliczenia $\pi_0$, mianowicie bijekcję $\pi_0^{-1}\circ\pi$ zbioru $\N_+$ na siebie. Z twierdzenia \ref{ma1:thm:absolute-rearrangement} każde takie przestawienie jest zbieżne bezwzględnie do tej samej sumy.
\end{proof}

\section{Wielokrotny iloczyn Cauchy'ego}
\begin{wniosek}
\label{ma1:cor:finite-cauchy-products}
Niech $r\in\N_+$ i niech dla każdego $i\in\{1,\ldots,r\}$ szereg $\sum_{n=0}^\infty a_n^{(i)}$ w $\R$ lub $\C$ będzie zbieżny bezwzględnie do $A_i$. Dla $n\ge0$ całkowitego połóżmy
\[
 d_n=\sum_{\substack{n_1,\ldots,n_r\in\N_0\\
                         n_1+\cdots+n_r=n}}
                    \prod_{i=1}^r a_{n_i}^{(i)}.
\]
Każda suma definiująca $d_n$ jest skończona, ponieważ wszystkie $n_i$ należą do $\{0,\ldots,n\}$. Wówczas szereg $\sum_{n=0}^\infty d_n$ jest zbieżny bezwzględnie i ma sumę $\prod_{i=1}^r A_i$. Ponadto
\[
 \sum_{n=0}^\infty |d_n|
 \le \prod_{i=1}^r\left(\sum_{n=0}^\infty |a_n^{(i)}|\right).
\]
\end{wniosek}
\begin{proof}
Dla $r=1$ teza jest bezpośrednia. Załóżmy ją dla $r$ szeregów i rozważmy $r+1$ szeregów. Oznaczmy współczynniki iloczynu pierwszych $r$ przez $d_j$. Współczynnik iloczynu Cauchy'ego tego szeregu i ostatniego szeregu jest równy
\[
 \sum_{j=0}^n d_j a_{n-j}^{(r+1)}
 =\sum_{\substack{n_1,\ldots,n_{r+1}\in\N_0\\
                        n_1+\cdots+n_{r+1}=n}}
             \prod_{i=1}^{r+1} a_{n_i}^{(i)}.
\]
Istotnie, każda krotka po prawej występuje dokładnie raz po lewej, dla $j=n_1+\cdots+n_r$. Wszystkie te sumy są skończone, bo każdy indeks należy do $\{0,\ldots,n\}$. Z założenia indukcyjnego szereg $\sum d_j$ jest bezwzględnie zbieżny. Stosując twierdzenie \ref{ma1:thm:absolute-cauchy-product}, otrzymujemy bezwzględną zbieżność nowego szeregu oraz iloczyn sum i oszacowanie sumy modułów przez iloczyn poprzedniego oszacowania i $\sum_{n\ge0}|a_n^{(r+1)}|$. Indukcja kończy dowód.
\end{proof}
\begin{uwaga}
Iloczyn Cauchy'ego nie oznacza mnożenia wyraz po wyrazie, czyli szeregu $\sum a_nb_n$. Pozwala on natomiast przejść od iloczynu dwóch szeregów do sumy współczynników wyznaczanych przez skończone sumy. Będzie to podstawą dowodu prawa dodawania argumentów funkcji wykładniczej. W następnym rozdziale zajmiemy się iloczynami nieskończonymi, definiowanymi jako granice ciągów iloczynów częściowych.
\end{uwaga}
